
\section{Control Strategies for Energy Management}

The development of effective control strategies for microgrid EMS has been a subject of extensive research. Broadly, these strategies can be classified into two main categories: optimization-based methods and heuristic rule-based approaches.

Optimization-based methods aim to find the mathematically optimal dispatch schedule by minimizing a cost objective function subject to system constraints. Linear Programming (LP) and Mixed-Integer Linear Programming (MILP) are among the most widely used techniques due to their ability to guarantee global optimality for convex problems. For instance, Sarda et al. \cite{SARDA2022280} demonstrated the effectiveness of LP in minimizing daily operational costs for a grid-tied microgrid. Similarly, Huang et al. \cite{Huang2021} applied MILP to optimize the operation of behind-the-meter storage, accounting for complex pricing structures. However, the reliance of these deterministic methods on perfect foresight is a major limitation. In real-time applications, accurate forecasts of load and generation are rarely available, and even minor prediction errors can lead to significant deviations from the optimal trajectory. Furthermore, the computational burden of these algorithms scales exponentially with the complexity of the system and the length of the prediction horizon, potentially rendering them unsuitable for rapid, real-time control loops.

In contrast, heuristic or rule-based strategies rely on pre-defined "if-then" logic derived from expert knowledge. A common example is the "self-consumption" strategy, where the battery charges whenever solar generation exceeds load and discharges when demand is high. Palma-Behnke et al. \cite{Palma2013} proposed a fuzzy logic controller that adjusts battery operation based on the state of charge and electricity price. These methods are computationally efficient and robust to forecast errors since they react to the current state of the system rather than relying on future predictions. However, their rigidity is a significant drawback. Rule-based controllers cannot anticipate future events, such as a spike in electricity prices or a period of cloud cover later in the afternoon. Consequently, they often fail to capture inter-temporal arbitrage opportunities, resulting in consistently suboptimal economic performance compared to optimization-based approaches \cite{Abushnaf2020}.

The limitations of both traditional paradigms have spurred interest in data-driven control techniques that can combine the optimality of mathematical programming with the real-time adaptability of heuristics.
